
\section{JUSTIFICACIÓN}

El sistema se justifica desde cuatro perspectivas: la social, la práctica, la técnica y la económica.

\subsection{Justificación social}

El Desafío Bebras acerca el pensamiento computacional a estudiantes de todo el país sin exigir conocimientos previos. Un sistema propio permite que cualquier unidad educativa participe con la guía de su maestro y que los estudiantes practiquen libremente antes y después de la prueba. Los formatos imprimibles llevan esa posibilidad a las aulas sin conexión.

\subsection{Justificación práctica}

Realizar el Desafío involucra procesos coordinados, desde la difusión de la convocatoria y el registro de las unidades educativas hasta la rendición y la publicación de resultados. El sistema reúne en un solo lugar la administración del sitio informativo, el registro de participantes, las categorías por edad, la rendición y el seguimiento de los resultados, y así da al equipo organizador los medios para administrar el Desafío por su cuenta.

\subsection{Justificación técnica}

Los sistemas de otros países tienen flujos funcionales válidos, pero, como muestra el análisis comparativo del capítulo 2, sus implementaciones responden a sus propias categorías, idiomas y formas de organizar la participación escolar, lo que desaconseja reutilizarlas directamente. El sistema desarrollado usa las seis categorías nacionales, el catálogo de unidades educativas del país y un sitio que el equipo organizador actualiza sin programar.

\subsection{Justificación económica}

El trabajo no demandó financiamiento. El sistema se construyó con herramientas de código abierto, sin costo de licencias, y la infraestructura para su despliegue (servidor, dominio y alojamiento) la provee la unidad patrocinante. La plataforma del Desafío funciona además en servicios en la nube, sin un servidor propio que comprar ni mantener.
